\chapter{Metodologia de trabajo}

La metodologia combina analisis del enunciado, entrevista simulada, extraccion de requisitos, diseño tecnico y evaluacion economica. El resultado es una propuesta de automatizacion teorica, pero redactada como si fuera una oferta preliminar para una planta industrial real.

\section{Pasos aplicados}

\begin{enumerate}
  \item \textbf{Revision del enunciado.} Se identificaron los entregables obligatorios: PDF profesional, presentacion, entrevista, oferta tecnica, oferta comercial, robot seleccionado, sensores adicionales y presupuesto detallado.
  \item \textbf{Entrevista simulada.} Se construyo una conversacion con Carlos Pérez, trabajador con mas de 10 años de experiencia en Alestis Puerto Real, para obtener requisitos operativos plausibles.
  \item \textbf{Extraccion de requisitos.} Los hallazgos se tradujeron en requisitos de navegacion, seguridad, maniobrabilidad, trazabilidad, inspeccion FOD, integracion y escalabilidad.
  \item \textbf{Diseño de solucion.} Se definio una arquitectura AMR omnidireccional con sensores de percepcion, trazabilidad, supervision y comunicacion industrial.
  \item \textbf{Evaluacion de despliegue.} Se plantearon tres escenarios: 1, 7 y 17 robots, con riesgos, beneficios y criterios de exito.
  \item \textbf{Oferta comercial y ROI.} Se elaboro un presupuesto por partidas y una estimacion prudente de recuperacion de tiempo operativo.
\end{enumerate}

\section{Fuentes y supuestos}

El enunciado de la Actividad 1 fija la entrega el 26 de mayo de 2026 antes de las 23:59, una ponderacion del 30\% en la asignatura y la necesidad de presentar oferta tecnica y comercial. Los datos economicos usados son los indicados por la actividad: produccion actual aproximada de 40 HTP al mes, objetivo futuro de 80 HTP al mes o mas, y salario medio de operario entre 20.000 y 28.000 \texteuro{} anuales.

Los importes del presupuesto son estimaciones academicas. Se han elegido valores conservadores para una solucion industrial con sensores, software, ingenieria, instalacion, formacion y mantenimiento, sin presentarlos como precios de proveedor.
